\subsection{Bird's-Eye View LiDAR Representation and Alternative Paradigms}
\label{sec:rel_bev_encoding}

Dense horizontal projection enables 2D convolutional networks to process unstructured 3D LiDAR point clouds. PIXOR~\cite{yang2018pixor} introduced proposal-free Bird's-Eye View (BEV) detection using fixed vertical occupancy slices and calibrated reflectance maps. Nguyen et al.~\cite{nguyen2024indoor} developed MobilePIXOR for indoor robotic perception, replacing large convolutional layers with inverted residual blocks and Gaussian heatmap supervision. PointPillars~\cite{lang2019pointpillars} instead organized points into vertical columns (pillars) and learned spatial features through local PointNet encoders. PillarNeXt~\cite{li2023pillarnext} and PillarNeSt~\cite{mao2023pillarnest} extended pillar-based receptive field design and backbone scaling for outdoor perception. Beyond Cartesian BEV representations, Range-View approaches such as LaserNet~\cite{meyer2019lasernet} represent LiDAR returns as spherical range images, enabling rapid processing with standard 2D convolutions. However, spherical projections suffer from severe metric scale distortion with radial distance and foreground-background occlusion tears. Conversely, modern query-based BEV Transformers such as TransFusion~\cite{bai2022transfusion} achieve state-of-the-art context modeling via sparse object queries cross-attending into dense LiDAR and camera feature maps. While query-to-feature cross-attention scales linearly with spatial token count ($O(N_{\mathrm{query}} \times HW)$), multi-head attention projections, high memory bandwidth, multi-sensor calibration synchronization, and large parameter footprints (frequently $20\text{--}40\,\text{M}$ parameters) far outstrip the strict sub-million-parameter limits of low-power edge microcontrollers.

BEV representations commonly use either learned PointNet-based pillar encoders or fixed-grid rasterization. Pillar encoders~\cite{lang2019pointpillars,li2023pillarnext} learn geometric features but require point gathering, dynamic memory allocation, and sparse scattering, which add latency and energy costs on edge accelerators. Fixed-grid methods~\cite{yang2018pixor,nguyen2024indoor} avoid learned point encoding. However, binary flags or raw cell counts on $0.1\,\text{m}$ grids omit continuous subcell offsets and higher-order vertical distributions, causing boundary staircasing and localization jitter. PillarHist~\cite{zhou2025pillarhist} investigated quantization-aware height histograms, and TriBand-BEV~\cite{khoshkdahan2026triband} explored multi-band vertical slicing for pedestrian detection. These methods do not combine continuous subcell centroid displacements with continuous vertical statistical moments in a compact tensor. MobilePIXORNeXt uses Strata-BEV14 (Section~\ref{sec:strata_bev14}) to combine four stratified height-band density bins ($F_0$--$F_3$), four continuous elevation moments ($z_{\max}, z_{\mathrm{mean}}, z_{\mathrm{span}}, z_{\mathrm{std}}$, $F_4$--$F_7$), two surface reflectance statistics ($I_{\max}, I_{\mathrm{mean}}$, $F_8$--$F_9$), global logarithmic cell density ($F_{10}$), binary occupancy ($F_{11}$), and two continuous subcell centroid displacement offsets $(\Delta\bar{x}, \Delta\bar{y}$, $F_{12}$--$F_{13}$). This 14-channel layout preserves sub-voxel geometry without learned point encoders, while reducing Host-to-Device transfer payloads by $60.0\%$ compared to 35-channel binary grids.

\subsection{Ground-Truth Data Augmentation and Physical Regularization}
\label{sec:rel_gt_augmentation}

To eliminate unphysical synthesis artifacts in LiDAR point-cloud augmentation, MobilePIXORNeXt introduces Physics-Consistent Ground-Truth Augmentation (\texttt{PC-GTA}, Section~\ref{sec:pc_gta}), integrating adaptive local ground plane fitting, static obstacle clearance filtering, line-of-sight ray checks, and shadow masking. While the ground-truth database sampling introduced in SECOND~\cite{yan2018second} enriches underrepresented pedestrian and cyclist classes, naive cut-and-paste insertion is unconstrained by environmental physics: sampled objects frequently float above ground, sink below road surfaces, or penetrate static background obstacles. Additionally, naive insertion violates raycasting optical geometry by failing to mask occluded background points behind newly inserted foreground targets and omitting distance-dependent beam divergence ($1/r^2$). By enforcing geometric and optical consistency, \texttt{PC-GTA} prevents these placement artifacts and contributes a $+6.03\,\text{pp}$ Moderate mAP improvement in controlled ablations.

\subsection{Lightweight Convolutional Backbones and Efficient Spatial Attention}
\label{sec:rel_backbone_attention}

Spatial feature extraction from sparse BEV planes must balance receptive field coverage against parameter and latency limits. For 3D LiDAR perception, SECOND~\cite{yan2018second} used sparse convolutions to process non-empty voxels. PointRCNN~\cite{shi2019pointrcnn}, Part-$A^2$~\cite{shi2020part}, and PV-RCNN~\cite{shi2020pvrcnn} developed two-stage point processing, part-aware aggregation, and voxel-set abstractions. These methods achieve high accuracy on KITTI~\cite{geiger2012kitti}, but their large parameter counts ($4.0\text{--}63.8\,\text{M}$) and substantial pipeline latencies ($93.7\text{--}194.5\,\text{ms}$) restrict real-time deployment on power-constrained edge accelerators.

Lightweight 2D convolutions offer an alternative to volumetric processing. ConvNeXt~\cite{liu2022convnext} modernized standard convolutional hierarchies using $7 \times 7$ depthwise kernels, inverted bottleneck expansions ($e=4$), GELU activations, and LayerScale, demonstrating accuracy competitive with Vision Transformers. For compact LiDAR BEV feature extraction, we tailor these architectural principles by adopting a slimmer expansion ratio ($e=2.5$), SiLU activations, and Batch Normalization, retaining large receptive fields while avoiding over-parameterization on edge hardware. Lightweight attention mechanisms further modulate geometric features. Efficient Channel Attention (ECA-Net)~\cite{wang2020eca} modeled 1D cross-channel interactions with low complexity but did not modulate individual spatial locations. Yang et al.~\cite{yang2021simam} proposed SimAM, deriving 3D spatial-channel attention analytically from an energy function inspired by visual neuroscience, without trainable parameters.

Prior mobile LiDAR detectors commonly use inverted residual bottlenecks with small $3 \times 3$ kernels~\cite{nguyen2024indoor}, which limit context across vehicle and pedestrian scales. External attention modules or distillation at coarse downsampling boundaries~\cite{tran2026rubber} add overhead while leaving receptive field limits within blocks unresolved under sparse point distributions. MobilePIXORNeXt combines the MobilePIXORNeXt-S backbone with parameter-free Block-SimAM (Sections~\ref{sec:backbone_architecture}--\ref{sec:next_s_block}) and a Top-Down Scale-Gated FPN (Section~\ref{sec:neck_architecture}). Large $7 \times 7$ depthwise kernels expand receptive fields, and attention modulates feature updates within each residual block without additional attention parameters.

\subsection{Geometric Supervision, Loss Formulations, and Quality Calibration}
\label{sec:rel_supervision_losses}

Oriented bounding box regression has developed from independent coordinate losses toward geometric objectives. Early detectors used decoupled Smooth-L1 losses. Independent regression of coordinates, dimensions, and angles ignores their geometric dependence and can destabilize gradients. Rezatofighi et al.~\cite{rezatofighi2019giou} introduced Generalized Intersection over Union (GIoU) to provide scale-invariant feedback when boxes do not overlap. Le et al.~\cite{le2026mgiou} later formulated Marginalized Generalized Intersection over Union (MGIoU), a projection-based objective for convex parametric shapes and arbitrarily oriented boxes without rotational singularities.

Localization quality calibration also affects detection ranking. Li et al.~\cite{li2020gfl} introduced Generalized Focal Loss (GFL). They showed that classification scores can disagree with localization quality, allowing high-confidence but poorly localized boxes to suppress accurate candidates during Non-Maximum Suppression (NMS). Kendall et al.~\cite{kendall2018uncertainty} pioneered multi-task uncertainty weighting by modeling homoscedastic observation noise, using a $\log \sigma$ penalty to prevent variance divergence.

Three supervision limitations remain in oriented detection. Independent angle and corner losses are sensitive to the $\pi$-periodicity boundary: small angular deviations near $\pm \pi/2$ or $\pm \pi$ can cause loss spikes and unstable yaw gradients. Auxiliary quality heads add parameters and can increase gradient competition across classes. In addition, unconstrained multi-task weighting can become unstable when classification, overlap, and orientation losses operate on disparate numerical scales. Unlike probabilistic observation noise formulations that require $\log \sigma$ regularization, MobilePIXORNeXt balances multi-task objectives via Temperature-Softmax Bounded Weighting (T-SBUW, Section~\ref{sec:losses_and_oga}). T-SBUW constrains task weights to a normalized, conserved simplex ($\sum \omega_k = J$) and employs running Exponential Moving Average (EMA) loss scale calibration to prevent weights from trivially collapsing toward tasks with smaller raw loss magnitudes. Together with Oriented Geometric Alignment (OGA) and IoU-Aware Quality Alignment (IQA), this framework provides numerically stable and balanced multi-task supervision.
